% supplementary_material_cosmology.tex
%
% Build with:  python3 paper/make_cosmology_figures.py && pdflatex supplementary_material_cosmology.tex
%
% Every number quoted below is written to figures/cosmology_facts.tex by paper/make_cosmology_figures.py,
% which runs the pyCICY modules, so the prose and the computations cannot drift apart silently.

\documentclass[11pt,a4paper]{article}
\usepackage[T1]{fontenc}
\usepackage[utf8]{inputenc}
\usepackage[margin=1in]{geometry}
\usepackage{amsmath,amssymb}
\usepackage{graphicx}
\usepackage{booktabs}
\usepackage[hidelinks]{hyperref}
\graphicspath{{figures/}}

% fallbacks so the document typesets before the facts exist; the \input overrides them
\providecommand{\CFPancakeContactTime}{?}\providecommand{\CFPancakePiBOverTwo}{?}
\providecommand{\CFPancakeFibreTurns}{?}\providecommand{\CFPancakeTin}{?}
\providecommand{\CFMixmasterEpochs}{?}\providecommand{\CFMixmasterLogHFinal}{?}
\providecommand{\CFMixmasterLogHChange}{?}\providecommand{\CFKasnerMaxP}{?}\providecommand{\CFKasnerMinP}{?}
\providecommand{\CFModuliSpeedDrift}{?}\providecommand{\CFELExponentDev}{?}\providecommand{\CFELConstraint}{?}
\providecommand{\CFDilatonELSigmaDev}{?}\providecommand{\CFDilatonELSpeedDev}{?}
\providecommand{\CFRacetrackSzero}{?}\providecommand{\CFRacetrackAlphaDown}{?}\providecommand{\CFRacetrackAlphaUp}{?}
\providecommand{\CFRacetrackLambdaDown}{?}\providecommand{\CFRacetrackLambdaUp}{?}
\providecommand{\CFDilationNorm}{?}\providecommand{\CFQuinticMetricAtTwo}{?}
\IfFileExists{figures/cosmology_facts.tex}{% generated by paper/make_cosmology_figures.py -- do not edit
\providecommand{\CFDilationNorm}{}\renewcommand{\CFDilationNorm}{1.5}
\providecommand{\CFELConstraint}{}\renewcommand{\CFELConstraint}{5.4e-18}
\providecommand{\CFELExponentDev}{}\renewcommand{\CFELExponentDev}{6.2e-14}
\providecommand{\CFKasnerMaxP}{}\renewcommand{\CFKasnerMaxP}{0.9767}
\providecommand{\CFKasnerMinP}{}\renewcommand{\CFKasnerMinP}{-0.3333}
\providecommand{\CFMixmasterEpochs}{}\renewcommand{\CFMixmasterEpochs}{40}
\providecommand{\CFMixmasterLogHChange}{}\renewcommand{\CFMixmasterLogHChange}{0.00e+00}
\providecommand{\CFMixmasterLogHFinal}{}\renewcommand{\CFMixmasterLogHFinal}{-2.622}
\providecommand{\CFModuliSpeedDrift}{}\renewcommand{\CFModuliSpeedDrift}{3.2e-12}
\providecommand{\CFPancakeContactTime}{}\renewcommand{\CFPancakeContactTime}{5.134}
\providecommand{\CFPancakeFibreTurns}{}\renewcommand{\CFPancakeFibreTurns}{1.464}
\providecommand{\CFPancakePiBOverTwo}{}\renewcommand{\CFPancakePiBOverTwo}{4.712}
\providecommand{\CFPancakeTin}{}\renewcommand{\CFPancakeTin}{1.0e-04}
\providecommand{\CFQuinticMetricAtTwo}{}\renewcommand{\CFQuinticMetricAtTwo}{0.375}
}{}

\title{Supplement: cosmological singularities and kinetic-dominated K\"ahler moduli in \texttt{pyCICY}}
\author{}
\date{}

\begin{document}
\maketitle

\section{What is here and what is not}
The modules \texttt{kasner}, \texttt{bkl}, \texttt{pancake}, \texttt{horizons}, \texttt{tidal}, \texttt{flrw} and
\texttt{ccbh} were ported from Quantum Gravity Lab (CC0). The sensor and ``complexity scalar'' of that project are not
ported: they map $K(t)$ to a number through a modelling choice, which is an experiment and not an exact result. The
toy cosmologies are our own, are not fitted, and are not reproductions of the models that inspired them. The one new
module, \texttt{moduli\_kasner}, is the only one that uses a Calabi--Yau: it takes the triple intersection numbers
$d_{ijk}$ and returns the exponent budget of section~\ref{sec:moduli}.

\section{Horizons}
\begin{figure}[h]\centering
\includegraphics[width=\textwidth]{cosmo_horizons.pdf}
\caption{Left: comoving null angle of the pancake along the collapsing fibre and along the base
($a_0=1,b_0=3,f_0=2,m=1/2$). Middle: fibre circumference. Right: cumulative Mixmaster horizon on each axis through
\CFMixmasterEpochs{} epochs of the sharp-wall BKL map, $u_0=\pi$ (irrational, since rational $u$ terminates the map).}
\end{figure}
The Kasner rule is exact: $\int_0^t t'^{-p}dt'$ is finite for $p<1$ and diverges logarithmically at $p=1$, so a
vacuum Kasner with $u>1$ has finite horizons on every axis (largest exponent for $u\in[1,6]$ is \CFKasnerMaxP{}, smallest
\CFKasnerMinP{}). The pancake's \emph{open} directions keep a finite horizon, reaching the antipodal angle $\pi$ of the
base at $t=\CFPancakeContactTime$ (compare $\pi b_0/2=\CFPancakePiBOverTwo$); the collapsing fibre is the one that
diverges, as $(2/a_0)\ln(t/t_{\rm in})$ with $t_{\rm in}=\CFPancakeTin$, a ray winding \CFPancakeFibreTurns{} times by $t=1$.
For Mixmaster the criterion $\int dt/a_i$ is finite on all three axes: the cumulative horizon changes by
\CFMixmasterLogHChange{} in $\log_{10}$ between the last two epochs, and the largest axis sits at
$\log_{10}h=\CFMixmasterLogHFinal$. That is a proxy criterion for Bianchi~IX, not a statement about causal contact in
the full solution.

\section{Moduli-driven Kasner exponents}\label{sec:moduli}
\paragraph{Setting.} Units $M_P=1$, action $R/2-K_{i\bar j}\partial T^i\partial\bar T^j$, classical large-volume
$K=-\ln(d_{ijk}t^it^jt^k)$, axions fixed (the slice $\mathrm{Im}\,T=0$ is totally geodesic, so the truncation is
consistent), no potential, Bianchi~I. In the real coordinates $t^i$
\begin{equation}
G_{ij}=\tfrac12\partial_i\partial_jK=-\frac{3\kappa_{ij}}{V}+\frac{9\kappa_i\kappa_j}{2V^2},\quad
V=d_{ijk}t^it^jt^k,\ \kappa_{ij}=d_{ijk}t^k,\ \kappa_i=d_{ijk}t^jt^k .
\end{equation}
\paragraph{Derivation.} (1) The shear equation is $\frac{d}{dt}(VH_i)=0$ with $V=a_1a_2a_3$, so
$H_i=p_i/(t-t_s)$ with $\sum p_i=1$. (2) The Hamiltonian constraint becomes
$\sum p_i^2=1-|s|^2$, where $s^a=d\phi^a/d\ln(t-t_s)$ and $|s|^2=G_{ab}s^as^b$ is conserved along the geodesic that the
moduli follow. (3) For $\sum p=1$ one has $\sum p^2\ge 1/3$, so real exponents exist iff
\begin{equation}
0\le |s|^2\le \tfrac23 ,
\end{equation}
with equality only at $p=(\tfrac13,\tfrac13,\tfrac13)$, the isotropic stiff fluid; at fixed $|s|$ the exponents lie on a
circle of squared radius $\tfrac23-|s|^2$ about that point. (4) The dilation $\xi=t^i\partial_i$ is Killing with
$G(\xi,\xi)=\tfrac32$ (computed as \CFDilationNorm{} at a test point), so its integral curves are geodesics and for
$\lambda\sim(t-t_s)^\sigma$ one has $|s|^2=\tfrac32\sigma^2$, i.e.\ $|\sigma|\le\tfrac23$: a kinetic-dominated
Calabi--Yau cannot grow its linear size faster than $(t-t_s)^{2/3}$. For the quintic ($h^{1,1}=1$, $d_{111}=5$) the
metric is $G_{11}=\tfrac{3}{2t^2}$, giving \CFQuinticMetricAtTwo{} at $t=2$.
\begin{figure}[h]\centering
\includegraphics[width=0.85\textwidth]{cosmo_moduli.pdf}
\caption{Left: circles of Kasner exponents for several $|s|^2$. Right: a numerical geodesic of the bicubic
($h^{1,1}=2$); $|s|^2$ drifts by \CFModuliSpeedDrift{} (relative) over $u\in[0,6]$.}
\end{figure}
\paragraph{Independent check.} The Euler--Lagrange equations of the Bianchi~I minisuperspace Lagrangian
$L=\tfrac12e^{S}[\sum\dot x_i^2-(\sum\dot x_i)^2+G_{ab}\dot t^a\dot t^b]$ are derived by sympy from $L$ alone and
integrated; the exponents are conserved to \CFELExponentDev{} and the constraint stays at \CFELConstraint.

\paragraph{The heterotic dilaton.} With $K\supset-\ln(S+\bar S)$ and $s=\mathrm{Re}\,S$ the dilaton factor is
$G_{ss}=1/(2s^2)$, flat in $\ln s$, and the total space is a Riemannian product, so the Kähler speed and
$\sigma_S=d\ln s/du$ are separately conserved and $|s|^2=|s_K|^2+\sigma_S^2/2\le\tfrac23$. The dilaton alone is bounded
by $\sigma_S^2\le\tfrac43$, and the joint dilation (norm $\tfrac32+\tfrac12=2$) by $\sigma^2\le\tfrac13$. The same
Euler--Lagrange integration with the dilaton added conserves $\sigma_S$ to \CFDilatonELSigmaDev{} and $|s|^2$ to
\CFDilatonELSpeedDev. Since $\alpha_{\rm GUT}=1/(4\pi\,\mathrm{Re}\,S)$, a dilaton growing by $F$ changes
$\alpha_{\rm GUT}$ by $F^{-\sigma_S}$, and at one loop $\Lambda_{\rm QCD}\propto\exp(-8\pi^2 s/|b_3|)$ exactly, matching
\texttt{running.lambda\_qcd}. As a thought experiment only, from the racetrack value $\mathrm{Re}\,S=\CFRacetrackSzero$
($N_1,N_2=4,5$, condensation ratio $10$ as input), using the whole budget for a factor $F=2$ would change
$\alpha_{\rm GUT}$ by $\CFRacetrackAlphaUp$ or $\CFRacetrackAlphaDown$ and $\Lambda_{\rm QCD}$ by $\CFRacetrackLambdaUp$ or
$\CFRacetrackLambdaDown$. The racetrack pins $\sigma_S=0$, so this bounds what a failure of stabilisation could do and
says nothing about whether it happens.

\paragraph{What is not claimed.} Geodesics for $h^{1,1}>1$ are numerical. The truncation omits $\alpha'$ corrections,
complex-structure and bundle moduli, axions and any potential; a heterotic dilaton enters only through the product
rule $|s|^2\to|s|^2+\tfrac12(d\ln s/du)^2$. Whether kinetic domination is reached at all is a question of moduli
stabilisation, which this package does not address, and nothing here is compared with data.

\end{document}
